\section{Proofs and derivations}
\label{app:proofs}

Notation follows Section~\ref{sec:identification}: $P_i$ and $q_i$ are distributions on $\{0,\dots,K-1\}$ with means $t_i$ and $m_i$, $F_q(k)=\sum_{r\le k}q(r)$ for $k=0,\dots,K-2$, and $d_k=F_P(k)-F_q(k)$.
Logarithms are natural; entropy $H$, cross-entropy $\CE$ and $\KL$ are in nats.
The statements are those of the identification preprint~\citep{miranda2026correlation}; we repeat the proofs so that the memo is self-contained.

\subsection{Pearson depends on means up to a positive affine map (Proposition~\ref{prop:affine})}
Since $\sum_i(t_i-\bar t)=0$, $\sum_i(m_i-\bar m)(t_i-\bar t)=\sum_i m_i(t_i-\bar t)=\sum_i(t_i-\bar t)\sum_r r\,q_i(r)$.
The numerator is linear in the probabilities and the denominator depends on the forecast only through $m$, so forecasts with equal item means have equal correlation.
The centered vector of $a\mathbf 1+bm$ is $b(m-\bar m\mathbf 1)$; numerator and standard deviation of the forecast both acquire the factor $b>0$, so $\rho$ is unchanged.
Cauchy--Schwarz bounds $\rho$ by one, with positive equality exactly when $m-\bar m\mathbf 1=b(t-\bar t\mathbf 1)$ for some $b>0$; setting $a=\bar m-b\bar t$ gives $m=a\mathbf 1+bt$.
The map $x\mapsto m(x)$ is unrestricted: the statement is about the output vector of means.

\subsection{Dimension of the unidentified set (Proposition~\ref{prop:dimension})}
For $K=2$, $q(1)=m$ and $q(0)=1-m$.
For general $K$ the constraints $\sum_rq(r)=1$ and $\sum_rr\,q(r)=m$ have linearly independent coefficient vectors, so their common affine solution set has dimension $K-2$; every interior mean $m\in(0,K-1)$ has an all-positive solution (the exponential family of \Cref{app:ce-gap}), so positivity adds no local equality constraint.
For $n$ items with prescribed interior means the dimensions add to $n(K-2)$.
In $m=a\mathbf 1+bt$ the columns $\mathbf 1$ and $t$ are independent because $t$ is nonconstant, which adds the two free parameters $(a,b)$ on the all-positive stratum, giving $n(K-2)+2$.
For binary ratings and correlation-one means $m_i=a+bt_i$, the Brier convention gives excess RPS $n^{-1}\sum_i(a+(b-1)t_i)^2$ and excess CE $n^{-1}\sum_i\KL(\mathrm{Bern}(t_i)\|\mathrm{Bern}(a+bt_i))$, so binary ratings keep level and scale ambiguity (the squared-error-over-both-categories convention would double the first expression).
For $K=3$ and mean one, the family $q_s=(s,1-2s,s)$, $0\le s\le\tfrac12$, has variance $2s$ at fixed mean.

\subsection{What each objective identifies (Proposition~\ref{prop:losses})}
\emph{Squared error.} For any $a$, $\E[(a-R)^2\mid x]=(a-\E[R\mid x])^2+\Var(R\mid x)$; the minimizer over $a$ is the conditional mean and the second term does not depend on $a$.

\emph{Cross-entropy.} $\CE(P,q)=H(P)+\KL(P\|q)$.
With $A=\{r:P(r)>0\}$ and $q$ positive on $A$, concavity of the logarithm gives $-\KL(P\|q)=\sum_{r\in A}P(r)\log(q(r)/P(r))\le\log\sum_{r\in A}q(r)\le0$, with equality only if $q/P$ is constant on $A$ and $q(A)=1$, that is $q=P$.
If $q$ vanishes on a category of positive $P$-mass, $\CE=+\infty$.

\emph{RPS.} Let $Z_k=\mathbf 1\{Y\le k\}$ with $Y\sim P$, so $\E Z_k=F_P(k)$ and $\Var Z_k=F_P(k)(1-F_P(k))$.
Because $\E(a-Z_k)^2=(a-F_P(k))^2+\Var Z_k$, summing over $k$ and dividing by $K-1$ gives
\[
\E_P\RPS(q,Y)=\frac1{K-1}\sum_kF_P(k)(1-F_P(k))+\frac{D(P,q)}{K-1},\qquad D=\sum_kd_k^2 .
\]
The first term does not depend on $q$, so the excess is $D/(K-1)$.
Since $\sum_{k=0}^{K-2}F_P(k)=\sum_rP(r)(K-1-r)=K-1-m_P$, we have $\sum_kd_k=m_q-m_P$.
The centered sum-of-squares identity $\sum_kd_k^2=(\sum_kd_k)^2/(K-1)+\sum_k(d_k-\bar d)^2$ then gives Equation~\eqref{eq:rps-decomp}.
The excess is zero exactly when all CDF values agree, which fixes every category probability by successive differences, so the normalized RPS is strictly proper on the whole simplex including its boundary.
Finally, by mass conservation $m_q-m_P=\sum_r(r-\tfrac{K-1}2)(q(r)-P(r))$ and so $|m_q-m_P|\le\tfrac{K-1}2\sum_r|q(r)-P(r)|=(K-1)\TVD(P,q)$.

\emph{TVD is not a proper score for individual ratings.}
For a point mass $\delta_y$, $\TVD(q,\delta_y)=1-q(y)$, so $\E_{Y\sim P}\TVD(q,\delta_Y)=1-\sum_rP(r)q(r)$, which is minimized by forecasts supported on the modal categories of $P$ and generally does not recover $P$.
TVD to a sparse empirical histogram is a noisy diagnostic whose expectation need not equal TVD to $P$.

\subsection{CE controls bounded moments (Proposition~\ref{prop:moments})}
Let $f:\{0,\dots,K-1\}\to[0,1]$, $\delta=P-q$ and $A=\{r:\delta(r)\ge0\}$.
Since $\sum_r\delta(r)=0$, $\sum_{r\in A}\delta(r)=\TVD(P,q)=d$, and $|\sum_rf(r)\delta(r)|\le d$.
For the second inequality assume $\KL(P\|q)<\infty$; grouping categories into $A$ and its complement cannot increase KL, by concavity of the logarithm.
With $a=P(A)$ and $b=q(A)$, so $a-b=d$, the grouped bound is the Bernoulli KL $h_b(a)$, with $h_b(b)=h_b'(b)=0$ and $h_b''(a)=1/(a(1-a))\ge4$; integrating twice gives $h_b(a)\ge2(a-b)^2$, so $\KL(P\|q)\ge2\TVD(P,q)^2$ (Pinsker).
Taking $f(r)=(r/(K-1))^j$, $|m_j(q,x)-m_j(p,x)|\le\TVD(p(\cdot|x),q(\cdot|x))\le\sqrt{\KL(p(\cdot|x)\|q(\cdot|x))/2}$, and averaging the square over $X$ with $\E_X\KL(p(\cdot|X)\|q(\cdot|X))=\CE(q)-\CE(p)$ gives the claim.

\subsection{Means alone cannot recover shape (Proposition~\ref{prop:lower})}
A rule that sees only the mean must output the same $q$ under $P_0$ and $P_1$.
The triangle inequality gives $d\le\TVD(P_0,q)+\TVD(q,P_1)$, so $\max_j\TVD(P_j,q)\ge d/2$.
Let $v_j=(F_{P_j}(k))_k$ and $u=(F_q(k))_k$.
The Euclidean identity $\tfrac12(\|v_0-u\|^2+\|v_1-u\|^2)=\|u-\tfrac12(v_0+v_1)\|^2+\tfrac14\|v_0-v_1\|^2$ and the RPS identity of the previous section give $\tfrac12(\mathcal E_{P_0}(q)+\mathcal E_{P_1}(q))\ge D(P_0,P_1)/(4(K-1))$.
For the $j$ attaining the maximum TVD, Pinsker gives $\KL(P_j\|q)\ge2(d/2)^2=d^2/2$.
The mixture $\tfrac12(P_0+P_1)$ attains the total variation bound and the average excess RPS bound when it is a feasible forecast.
These are information bounds for a rule that receives only the mean; they do not bound an estimator that also sees individual ratings, features or shape restrictions.

\subsection{Gradients}
Let $q_i=\mathrm{softmax}(u_i)$, $\zeta_i=(t_i-\bar t)/s_t$, $z_i=(m_i-\bar m)/s_m$, $T=t-\bar t\mathbf 1$, $C=m-\bar m\mathbf 1$ and $N=T^\top m$, so $\rho=N/(\|T\|\|C\|)$.
Then $\partial N/\partial m_i=T_i$ and $\partial\|C\|/\partial m_i=C_i/\|C\|$ because $\sum_iC_i=0$, and the quotient rule gives $\partial\rho/\partial m_i=g_i=(\zeta_i-\rho z_i)/(ns_m)$.
The softmax derivative $\partial q_i(s)/\partial u_i(r)=q_i(s)(\mathbf 1\{s=r\}-q_i(r))$ gives $\partial m_i/\partial u_i(r)=q_i(r)(r-m_i)$, hence $\partial\rho/\partial u_i(r)=g_iq_i(r)(r-m_i)$.
Differentiating $-\sum_sP_i(s)\log q_i(s)$ gives $q_i(r)-P_i(r)$.
For a differentiable shared-parameter model, $D\rho(\theta)[v]=\sum_ig_iDM_i(\theta)[v]$ vanishes whenever $DM(\theta)[v]=0$; this is a first-order statement at a point and not a statement about higher-order changes.

\subsection{The exponential family, the signed CE decomposition and the entropy gap}
\label{app:ce-gap}
Let $q_\gamma(r)=e^{\gamma r-A(\gamma)}$ on $\{0,\dots,K-1\}$ with $A(\gamma)=\log\sum_re^{\gamma r}$.
Then $A'(\gamma)=m(\gamma)$ and $A''(\gamma)=\Var_{q_\gamma}(R)>0$, so the mean map is a strictly increasing bijection from $\mathbb R$ onto $(0,K-1)$.
For a target $P$ with interior mean $t$, $\CE(P,q_\gamma)=-\gamma t+A(\gamma)$ is strictly convex in $\gamma$ with derivative $m(\gamma)-t$, so its unique minimizer is the mean-matching member $q_t$.
For any $P$ with mean $t$, $\KL(P\|q_t)=-H(P)-\gamma t+A(\gamma)=H(q_t)-H(P)$, so $q_t$ is the unique maximum-entropy distribution with mean $t$.
Define the gap $G(P)=H(q_t)-H(P)\ge0$: it is the excess CE of the best forecast that is log-affine in the rating.
For any $q$ positive on every category, adding and subtracting $\CE(q_t,q)=H(q_t)+\KL(q_t\|q)$ gives
\[
\CE(P,q)=H(P)+G(P)+\KL(q_t\|q)+C(P,q),\qquad C(P,q)=\sum_r(P(r)-q_t(r))(-\log q(r)).
\]
The cross term vanishes when $\log q(r)=\alpha+\beta r$, because $P$ and $q_t$ share total mass and mean, and it can be negative: at $q=P$, $C(P,P)=-G(P)-\KL(q_t\|P)$.
The gap is therefore a restriction cost for the log-affine family and not an unavoidable penalty on every forecast.
In the 11-artifact population of experiment 04 ($p_x=0.9\,\mathrm{Binomial}(5,j/12)+0.1\,\mathrm{Uniform}$), the mean gap is $0.117$ nats (item range $0.020$ to $0.210$), against the largest possible entropy $\log6=1.792$ nats.
Locally, $\KL(q_{\gamma_t}\|q_{\gamma_m})=(m-t)^2/(2A''(\gamma_t))+o((m-t)^2)$, the cost of a mean miscalibration inside the family.

\subsection{The exact example}
\label{app:example}
\Cref{tab:example} lists the three items of \Cref{fig:counterexample}; every quantity is exact rational arithmetic except CE.
\begin{table}[h]
\caption{\textbf{Identical means and Pearson correlation 1, different variances and distributions.}
Exact counterexample, $K=3$, $n=3$.
$P_i$: human distribution; $Q_i$: forecast; means and variances of the ratings 0, 1, 2; $\TVD_i$: total variation distance between $P_i$ and $Q_i$ (0--1); excess ranked probability score (RPS) of item $i$, written excess RPS$_i$: the expected RPS of $Q_i$ minus that of $P_i$ under $P_i$ (0--1).
Averages over the three items: mean TVD $1/5$; excess RPS $3/200$; true RPS $49/300$, forecast RPS $107/600$; cross-entropy (CE) $0.915387$ for $P$ and $1.064735$ nats for $Q$.}
\label{tab:example}
\centering\small
\begin{tabular}{llccccc}
\toprule
Item & $P_i$ & $Q_i$ & mean & variance $P_i$ / $Q_i$ & $\TVD_i$ & excess RPS$_i$ \\
\midrule
1 & $(3/5,\,3/10,\,1/10)$ & $(13/20,\,1/5,\,3/20)$ & $1/2$ & $9/20\ /\ 11/20$ & $1/10$ & $1/400$ \\
2 & $(1/5,\,3/5,\,1/5)$ & $(2/5,\,1/5,\,2/5)$ & $1$ & $2/5\ /\ 4/5$ & $2/5$ & $1/25$ \\
3 & $(1/10,\,3/10,\,3/5)$ & $(3/20,\,1/5,\,13/20)$ & $3/2$ & $9/20\ /\ 11/20$ & $1/10$ & $1/400$ \\
\bottomrule
\end{tabular}
\end{table}
Both mean vectors equal $(1/2,1,3/2)$, so $\rho=1$ by Proposition~\ref{prop:affine}.
In every item the mean term of Equation~\eqref{eq:rps-decomp} is zero, so the whole excess $(1+16+1)/(400\cdot3)=3/200$ is the second term.

\subsection{Finite ratings and the heuristic break-even}
\label{app:finite}
If each held-out item has $k$ conditionally independent ratings with mean $\hat t_I=t_I+\varepsilon_I$, $\E[\varepsilon_I\mid I]=0$ and $\Var(\varepsilon_I\mid I)=\sigma_I^2/k$, and the forecast mean $m_I$ is independent of the test-rating noise, then $\Cov(m_I,\hat t_I)=\Cov(m_I,t_I)$ and $\Var(\hat t_I)=\Var(t_I)+\E[\sigma_I^2]/k$.
The correlation with the observed mean is thus attenuated by $\sqrt{\Var(t_I)/(\Var(t_I)+\E[\sigma_I^2]/k)}$, which is why the same judge looks worse against one rating per item than against five.
For a well-specified parametric conditional model with $P$ identifiable parameters, an efficient consistent maximum-likelihood estimator whose scaled error $\sqrt N(\hat\theta-\theta^\ast)$ has converging second moments, nonsingular Fisher information and expected Taylor remainder $o(N^{-1})$, the expected excess CE is $\E[R(\hat\theta)-R(\theta^\ast)]=P/(2N)+o(N^{-1})$ nats per fresh observation.
If a specified procedure converges to a forecast with positive excess CE $B$, comparing $P/(2N)$ with $B$ suggests the heuristic scale $N\approx P/(2B)$.
It requires both risks in CE nats and the asymptotic assumptions to hold; it fails when $N$ is comparable to $P$, near a boundary, under misspecification or strong regularization, and an RPS gap cannot replace $B$.
Properness describes the optimizer of population loss, not the order of finite-sample estimators, so neither $B>0$ nor a universal crossover follows from the invariances of Proposition~\ref{prop:affine}.
